%
%============= CHAPTER BREAK =================%
%
\chapter{Appian of Alexandria (147--161 CE)}%
\vspace*{3mm}%
\justifying%
%
\noindent{}Appian of Alexandria wrote his massive \textit{Historia Romana} no earlier %
than 147 CE, after his appointment as procurator, %
and some 900 years after Romeʼs founding, as he claims.\endnote{%
Preface 6--7; White 1912:8--11. His chronological estimates subdivided into three major epochs (500+200+200) %
from Romeʼs founding until his composition. The self-reference as a procurator concludes %
the preface at section 15; White 1912:24--25.} %
This monumental work was likely completed no later than 162 CE, %
since it shows no awareness of later events.\endnote{%
E.g., the reintroduction by Marcus Aurelius in 162 CE of %
proconsuls ruling parts of Italy goes unmentioned in \textit{Roman History} 13.38. %
See Livius.org, s.v., “Appian”.} %
\vfill{}

\noindent{}\textbf{ἐπιπόθητος} / \textit{epipothētos} / “yearned for” yields 102 TLG hits. %
It is cognate to ἐπιπόθησις / \textit{epipothēsis} / “yearning”, covered in the chapter on Favorinus. %
Appian stands as the second non-Christian author to use the former lemma, %
but the first to use the lemma \textit{epipothētos}, %
specifically in the \textit{Roman History}, in the book \textit{Iberica}.\endnote{McGing 2019:204, 206.} %

\vfill{}
 \begin{quote}Φλάκκῳ μὲν οὖν διάδοχος ἦλθεν ἐπὶ τὴν στρατηγίαν Τιβέριος Σεμπρώνιος Γράκχος. ... %
 Τοὺς δὲ ἀπόρους συνῴκιζε καὶ γῆν αὐτοῖς διεμέτρει %
 καὶ πᾶσιν ἔθετο τοῖς τῇδε συνθήκας ἀκριβεῖς, %
 καθ’ ἃς Ῥωμαίων ἔσονται φίλοι· %
 ὅρκους τε ὤμοσεν αὐτοῖς καὶ ἔλαβεν, \textbf{ἐπιποθήτους} ἐν τοῖς ὕστερον πολέμοις πολλάκις γενομένους. %
 Δι’ ἃ καὶ ἐν Ἰβηρία καὶ ἐν Ῥώμῃ διώνυμος ἐγένετο ὁ Γράκχος καὶ ἐθριάμβευσε λαμπρῶς. %
 (\textit{H.R. Iber.} 43.175, 179) \\[1mm] %
 So then, Tiberius Sempronius Gracchus came as successor to Flaccus in military command. ... %
 He was uniting the unmanageable ones and was allotting them land and was setting strict treaties with all of them present, %
 according to which they would be friends of the Romans. %
 And he swore oaths to them and received [them in return], %
 [oaths] many times \textbf{yearned for} in later wars. %
 Based on these [accomplishments], %
 both in Iberia and in Rome, Gracchus became far-famed and brilliantly led a triumph-procession.\end{quote} % 
 \clearpage

\noindent{}The same, apparently new word occurs in canonical Philippians 4:1, %
in a passage not evidently present in Marcionʼs \textit{Apostolos}. %

\vspace*{3mm} 
 \begin{quote} Ὥστε, ἀδελφοί μου ἀγαπητοὶ καὶ \textbf{ἐπιπόθητοι}, %
    χαρὰ καὶ στέφανός μου, οὕτως στήκετε ἐν κυρίῳ, ἀγαπητοί. %
  ² Εὐοδίαν παρακαλῶ καὶ Συντύχην παρακαλῶ τὸ αὐτὸ φρονεῖν ἐν κυρίῳ. %
  ³ ναὶ ἐρωτῶ καὶ σέ, γνήσιε σύζυγε, συλλαμβάνου αὐταῖς, %
  αἵτινες ἐν τῷ εὐαγγελίῳ συνήθλησάν μοι μετὰ καὶ Κλήμεντος %
  καὶ τῶν λοιπῶν συνεργῶν μου, ὧν τὰ ὀνόματα ἐν βίβλῳ ζωῆς. (Phlp 4:1--3) \\[1mm] %
  So then, my beloved and \textbf{yearned for} brothers, joy and crown of mine, thus stand in Lord, brothers. %
  ² Euodia I urge and Syntyche I urge, to be minded the same in Lord. %
  ³ Yes, I ask you also, genuine yoke-fellow, take up together with them, %
  they who co-competed with me in the Gospel along with Clement and the other co-workers of mine, %
  whose names [are] in the book of life. \end{quote} % 
 \vspace*{3mm}

Note not only the parallel new lemmas, %
but also similar themes, including the unification of hostile parties and celebration of victory. %
Canonical Paul in effect becomes a legendary rhetorical general, %
bringing harmony to relationships in the Christian community %
the way Gracchus forged alliances in Iberia. %
Against this background, the name \textit{Syntyche} / “Shared-Fortune” %
also cleverly hints at the formation of a treaty after conflict, a cessation from hostilities. %  
The mention of a certain “Clement” in this context is also a crossover redactional tell. %
Notice the same new lemma appears in Ps-Clement, \textit{1Cor} 65, %
in a section that begins abruptly after a doxological conclusion in \textit{1Cor} 64, %
one that also contains a Fortune-related name. %

\vspace*{3mm} 
 \begin{quote}Τοὺς δὲ ἀπεσταλμένους ἀφ’ ἡμῶν Κλαύδιον Ἔφηβον καὶ Οὐαλέριον Βίτωνα %
 σὺν καὶ Φορτουνάτῳ ἐν εἰρήνῃ μετὰ χαρᾶς ἐν τάχει ἀναπέμψατε πρὸς ἡμᾶς, %
 ὅπως θᾶττον τὴν εὐκταίαν καὶ \textbf{ἐπιποθήτην} ἡμῖν εἰρήνην καὶ ὁμόνοιαν ἀπαγγέλλωσιν, %
 εἰς τὸ τάχιον καὶ ἡμᾶς χαρῆναι περὶ τῆς εὐσταθείας ὑμῶν. \\[1mm] %
 Now them sent forth by us --- Claudius Ephebus and Valerius Bito, %
 along with Fortunatus --- send back unto us in peace with joy in haste, %
 so that they may more hastily report unto us the votive and peace and harmony \textbf{yearned for} by us, %
 in order that we be made joyful more hastily concerning the tranquility of yours. \end{quote} % 
 \vspace*{3mm}

A secondary, yet again strictly canonical word-link deepens the parallel between %
the canonical Pauline redaction and Appianʼs Gracchus. %
Notice here the word εὐωδία / \textit{euōdia} / “fragrance” (identical to the personal name Euodia) %
and the verb “to lead a triumph-procession” / θριαμβεύω / \textit{thriambeuō}. %
Cp. θριαμβεύοντι / \textit{thriambeuonti} (2Co 2:14) and ἐθριάμβευσε / \textit{ethriambeuse} (\textit{H.R. Iber.} 43.179). %
The fuller canonical quotation follows. %
Note that this segment was missing from Marcionʼs \textit{Apostolos}, %
even though it picked up immediately after with 2Cor 2:15b. %

\vspace*{3mm} 
 \begin{quote}Τῷ δὲ θεῷ χάρις τῷ πάντοτε θριαμβεύοντι ἡμᾶς ἐν τῷ Χριστῷ %
 καὶ τὴν ὀσμὴν τῆς γνώσεως αὐτοῦ φανεροῦντι δι᾽ ἡμῶν ἐν παντὶ τόπῳ·
  ¹⁵ ὅτι Χριστοῦ εὐωδία ἐσμὲν τῷ θεῷ (2Cor 2:14--15) \\[1mm] %
 Now favor be to the God who always leads us in triumphal procession in the Christ %
 and through us in every place manifests the scent of the knowledge of him, 
  ¹⁵ for we are Christʼs fragrance to the God.\end{quote} % 
\vspace*{3mm}

ἐπιπόθητος / \textit{epipothētos} is also found in \textit{Ep.Barn.} 1.3b and the \textit{Acts of Andrew} 10. %
Both passages echo additional language and stylistic tendencies found in Phlp 4:1--3 and Ps-Clem, \textit{1Cor}, %
such as terms related to haste, peace, joy, and solidarity, %
as well as words gratuitously prefixed with συ-/\textit{su}-, translated “con-|with”. %